% Options for packages loaded elsewhere
% Options for packages loaded elsewhere
\PassOptionsToPackage{unicode}{hyperref}
\PassOptionsToPackage{hyphens}{url}
\PassOptionsToPackage{dvipsnames,svgnames,x11names}{xcolor}
%
\documentclass[
  letterpaper,
  DIV=11,
  numbers=noendperiod]{scrartcl}
\usepackage{xcolor}
\usepackage{amsmath,amssymb}
\setcounter{secnumdepth}{-\maxdimen} % remove section numbering
\usepackage{iftex}
\ifPDFTeX
  \usepackage[T1]{fontenc}
  \usepackage[utf8]{inputenc}
  \usepackage{textcomp} % provide euro and other symbols
\else % if luatex or xetex
  \usepackage{unicode-math} % this also loads fontspec
  \defaultfontfeatures{Scale=MatchLowercase}
  \defaultfontfeatures[\rmfamily]{Ligatures=TeX,Scale=1}
\fi
\usepackage{lmodern}
\ifPDFTeX\else
  % xetex/luatex font selection
\fi
% Use upquote if available, for straight quotes in verbatim environments
\IfFileExists{upquote.sty}{\usepackage{upquote}}{}
\IfFileExists{microtype.sty}{% use microtype if available
  \usepackage[]{microtype}
  \UseMicrotypeSet[protrusion]{basicmath} % disable protrusion for tt fonts
}{}
\makeatletter
\@ifundefined{KOMAClassName}{% if non-KOMA class
  \IfFileExists{parskip.sty}{%
    \usepackage{parskip}
  }{% else
    \setlength{\parindent}{0pt}
    \setlength{\parskip}{6pt plus 2pt minus 1pt}}
}{% if KOMA class
  \KOMAoptions{parskip=half}}
\makeatother
% Make \paragraph and \subparagraph free-standing
\makeatletter
\ifx\paragraph\undefined\else
  \let\oldparagraph\paragraph
  \renewcommand{\paragraph}{
    \@ifstar
      \xxxParagraphStar
      \xxxParagraphNoStar
  }
  \newcommand{\xxxParagraphStar}[1]{\oldparagraph*{#1}\mbox{}}
  \newcommand{\xxxParagraphNoStar}[1]{\oldparagraph{#1}\mbox{}}
\fi
\ifx\subparagraph\undefined\else
  \let\oldsubparagraph\subparagraph
  \renewcommand{\subparagraph}{
    \@ifstar
      \xxxSubParagraphStar
      \xxxSubParagraphNoStar
  }
  \newcommand{\xxxSubParagraphStar}[1]{\oldsubparagraph*{#1}\mbox{}}
  \newcommand{\xxxSubParagraphNoStar}[1]{\oldsubparagraph{#1}\mbox{}}
\fi
\makeatother

\usepackage{color}
\usepackage{fancyvrb}
\newcommand{\VerbBar}{|}
\newcommand{\VERB}{\Verb[commandchars=\\\{\}]}
\DefineVerbatimEnvironment{Highlighting}{Verbatim}{commandchars=\\\{\}}
% Add ',fontsize=\small' for more characters per line
\usepackage{framed}
\definecolor{shadecolor}{RGB}{241,243,245}
\newenvironment{Shaded}{\begin{snugshade}}{\end{snugshade}}
\newcommand{\AlertTok}[1]{\textcolor[rgb]{0.68,0.00,0.00}{#1}}
\newcommand{\AnnotationTok}[1]{\textcolor[rgb]{0.37,0.37,0.37}{#1}}
\newcommand{\AttributeTok}[1]{\textcolor[rgb]{0.40,0.45,0.13}{#1}}
\newcommand{\BaseNTok}[1]{\textcolor[rgb]{0.68,0.00,0.00}{#1}}
\newcommand{\BuiltInTok}[1]{\textcolor[rgb]{0.00,0.23,0.31}{#1}}
\newcommand{\CharTok}[1]{\textcolor[rgb]{0.13,0.47,0.30}{#1}}
\newcommand{\CommentTok}[1]{\textcolor[rgb]{0.37,0.37,0.37}{#1}}
\newcommand{\CommentVarTok}[1]{\textcolor[rgb]{0.37,0.37,0.37}{\textit{#1}}}
\newcommand{\ConstantTok}[1]{\textcolor[rgb]{0.56,0.35,0.01}{#1}}
\newcommand{\ControlFlowTok}[1]{\textcolor[rgb]{0.00,0.23,0.31}{\textbf{#1}}}
\newcommand{\DataTypeTok}[1]{\textcolor[rgb]{0.68,0.00,0.00}{#1}}
\newcommand{\DecValTok}[1]{\textcolor[rgb]{0.68,0.00,0.00}{#1}}
\newcommand{\DocumentationTok}[1]{\textcolor[rgb]{0.37,0.37,0.37}{\textit{#1}}}
\newcommand{\ErrorTok}[1]{\textcolor[rgb]{0.68,0.00,0.00}{#1}}
\newcommand{\ExtensionTok}[1]{\textcolor[rgb]{0.00,0.23,0.31}{#1}}
\newcommand{\FloatTok}[1]{\textcolor[rgb]{0.68,0.00,0.00}{#1}}
\newcommand{\FunctionTok}[1]{\textcolor[rgb]{0.28,0.35,0.67}{#1}}
\newcommand{\ImportTok}[1]{\textcolor[rgb]{0.00,0.46,0.62}{#1}}
\newcommand{\InformationTok}[1]{\textcolor[rgb]{0.37,0.37,0.37}{#1}}
\newcommand{\KeywordTok}[1]{\textcolor[rgb]{0.00,0.23,0.31}{\textbf{#1}}}
\newcommand{\NormalTok}[1]{\textcolor[rgb]{0.00,0.23,0.31}{#1}}
\newcommand{\OperatorTok}[1]{\textcolor[rgb]{0.37,0.37,0.37}{#1}}
\newcommand{\OtherTok}[1]{\textcolor[rgb]{0.00,0.23,0.31}{#1}}
\newcommand{\PreprocessorTok}[1]{\textcolor[rgb]{0.68,0.00,0.00}{#1}}
\newcommand{\RegionMarkerTok}[1]{\textcolor[rgb]{0.00,0.23,0.31}{#1}}
\newcommand{\SpecialCharTok}[1]{\textcolor[rgb]{0.37,0.37,0.37}{#1}}
\newcommand{\SpecialStringTok}[1]{\textcolor[rgb]{0.13,0.47,0.30}{#1}}
\newcommand{\StringTok}[1]{\textcolor[rgb]{0.13,0.47,0.30}{#1}}
\newcommand{\VariableTok}[1]{\textcolor[rgb]{0.07,0.07,0.07}{#1}}
\newcommand{\VerbatimStringTok}[1]{\textcolor[rgb]{0.13,0.47,0.30}{#1}}
\newcommand{\WarningTok}[1]{\textcolor[rgb]{0.37,0.37,0.37}{\textit{#1}}}

\usepackage{longtable,booktabs,array}
\usepackage{calc} % for calculating minipage widths
% Correct order of tables after \paragraph or \subparagraph
\usepackage{etoolbox}
\makeatletter
\patchcmd\longtable{\par}{\if@noskipsec\mbox{}\fi\par}{}{}
\makeatother
% Allow footnotes in longtable head/foot
\IfFileExists{footnotehyper.sty}{\usepackage{footnotehyper}}{\usepackage{footnote}}
\makesavenoteenv{longtable}
\usepackage{graphicx}
\makeatletter
\newsavebox\pandoc@box
\newcommand*\pandocbounded[1]{% scales image to fit in text height/width
  \sbox\pandoc@box{#1}%
  \Gscale@div\@tempa{\textheight}{\dimexpr\ht\pandoc@box+\dp\pandoc@box\relax}%
  \Gscale@div\@tempb{\linewidth}{\wd\pandoc@box}%
  \ifdim\@tempb\p@<\@tempa\p@\let\@tempa\@tempb\fi% select the smaller of both
  \ifdim\@tempa\p@<\p@\scalebox{\@tempa}{\usebox\pandoc@box}%
  \else\usebox{\pandoc@box}%
  \fi%
}
% Set default figure placement to htbp
\def\fps@figure{htbp}
\makeatother





\setlength{\emergencystretch}{3em} % prevent overfull lines

\providecommand{\tightlist}{%
  \setlength{\itemsep}{0pt}\setlength{\parskip}{0pt}}



 


\KOMAoption{captions}{tableheading}
\makeatletter
\@ifpackageloaded{caption}{}{\usepackage{caption}}
\AtBeginDocument{%
\ifdefined\contentsname
  \renewcommand*\contentsname{Table of contents}
\else
  \newcommand\contentsname{Table of contents}
\fi
\ifdefined\listfigurename
  \renewcommand*\listfigurename{List of Figures}
\else
  \newcommand\listfigurename{List of Figures}
\fi
\ifdefined\listtablename
  \renewcommand*\listtablename{List of Tables}
\else
  \newcommand\listtablename{List of Tables}
\fi
\ifdefined\figurename
  \renewcommand*\figurename{Figure}
\else
  \newcommand\figurename{Figure}
\fi
\ifdefined\tablename
  \renewcommand*\tablename{Table}
\else
  \newcommand\tablename{Table}
\fi
}
\@ifpackageloaded{float}{}{\usepackage{float}}
\floatstyle{ruled}
\@ifundefined{c@chapter}{\newfloat{codelisting}{h}{lop}}{\newfloat{codelisting}{h}{lop}[chapter]}
\floatname{codelisting}{Listing}
\newcommand*\listoflistings{\listof{codelisting}{List of Listings}}
\makeatother
\makeatletter
\makeatother
\makeatletter
\@ifpackageloaded{caption}{}{\usepackage{caption}}
\@ifpackageloaded{subcaption}{}{\usepackage{subcaption}}
\makeatother
\usepackage{bookmark}
\IfFileExists{xurl.sty}{\usepackage{xurl}}{} % add URL line breaks if available
\urlstyle{same}
\hypersetup{
  pdftitle={Lab: Hormone Supplements and Disease},
  colorlinks=true,
  linkcolor={blue},
  filecolor={Maroon},
  citecolor={Blue},
  urlcolor={Blue},
  pdfcreator={LaTeX via pandoc}}


\title{Lab: Hormone Supplements and Disease}
\usepackage{etoolbox}
\makeatletter
\providecommand{\subtitle}[1]{% add subtitle to \maketitle
  \apptocmd{\@title}{\par {\large #1 \par}}{}{}
}
\makeatother
\subtitle{Power and Pitfalls in Experiments}
\author{}
\date{}
\begin{document}
\maketitle


\subsection{Introduction}\label{introduction}

The goal of this exercise is to better understand the concept of
statistical power, and how easily it may be misconstrued in real world
experiments. Power is of paramount importance in a lot of applied
research, but despite its easy theoretical definition and guarantees,
there is a lot of room for bias and other errors to creep in. We shall
look at one such well-known case through the years, and read through and
analyze it ourselves before doing some simple coding with dummy
datasets.

\subsection{The Experiments}\label{the-experiments}

We will be looking at influential papers that attempt to analyze the
effect of post-menopausal hormone therapy in women with coronary and
cardiovascular disease risk, to better understand the effects of such
supplements. These were conducted between 1991 to 2007, though this
topic has been studied after this time period as well.

We first look at two studies: The Nurses' Health Study from 1991 and the
WHI (Women's Health Initiative) study from 2002. Read through both
papers and answer the following questions.

\begin{enumerate}
\def\labelenumi{\arabic{enumi}.}
\item
  Compare the four components of the experimental design between both
  studies. Does The Nurses Health Study qualify as an experiment? Do its
  assumptions allow for causal inference?

  \begin{itemize}
  \item
    The Nature of the Treatments.
  \item
    Choice of the Experimental Units.
  \item
    Manner of Assigning Units to Treatments
  \item
    The Nature of the Response
  \end{itemize}
\end{enumerate}

\hfill\break
\hfill\break
\hfill\break
\hfill\break
\hfill\break
\hfill\break
\hfill\break
\hfill\break
\hfill\break
\hfill\break

\begin{enumerate}
\def\labelenumi{\arabic{enumi}.}
\setcounter{enumi}{1}
\tightlist
\item
  Both studies analyze the same question, but arrive at very different
  conclusions. Where do you think this difference arises from?
\end{enumerate}

\hfill\break
\hfill\break
\hfill\break
\hfill\break
\hfill\break
\hfill\break
\hfill\break

\begin{enumerate}
\def\labelenumi{\arabic{enumi}.}
\setcounter{enumi}{2}
\tightlist
\item
  Which study do you trust more and why or why not? Where do you think
  the other study fails in its attempt to answer the question?
\end{enumerate}

\hfill\break
\hfill\break
\hfill\break
\hfill\break
\hfill\break
\hfill\break
\hfill\break
\hfill\break

\begin{enumerate}
\def\labelenumi{\arabic{enumi}.}
\setcounter{enumi}{3}
\tightlist
\item
  Which of these studies has a larger sample size? Which would you say
  is more powerful?
\end{enumerate}

\hfill\break
\hfill\break
\hfill\break
\hfill\break
\hfill\break
\hfill\break
\hfill\break

\begin{enumerate}
\def\labelenumi{\arabic{enumi}.}
\setcounter{enumi}{4}
\tightlist
\item
  As visible above, different applied research processes, despite both
  seeming statistically valid, can often lead to entirely opposite
  conclusions with statistical significance. Does the statistical theory
  of power then fail here?
\end{enumerate}

\hfill\break
\hfill\break
\hfill\break
\hfill\break
\hfill\break
\hfill\break
\hfill\break
\hfill\break

\subsection{A Third Study}\label{a-third-study}

\hfill\break
The 2002 results from the Women's Health Initiative (WHI) dramatically
changed how doctors viewed hormone therapy. Unlike earlier observational
research such as the Nurses' Health Study, which suggested hormones
might protect the heart, the WHI trial found increased health risks.

Thus, in 2007 researchers revisited the data from the WHI study to
explore whether subgroups by age could explain the disagreement, testing
the idea that hormone therapy might have different effects depending on
when it is started. Read the study and answer the following:

\begin{enumerate}
\def\labelenumi{\arabic{enumi}.}
\setcounter{enumi}{5}
\tightlist
\item
  Compare the overall hazard ratio reported in 2002 with the hazard
  ratios reported for the 50--59, 60--69, and 70--79 age groups in 2007.
  What differences do you observe?
\end{enumerate}

\hfill\break
\hfill\break
\hfill\break
\hfill\break
\hfill\break
\hfill\break
\hfill\break
\hfill\break

\begin{enumerate}
\def\labelenumi{\arabic{enumi}.}
\setcounter{enumi}{6}
\tightlist
\item
  What questions do you think the two studies try to answer? Are they
  different? How does this change their power, and the validity of their
  conclusions?
\end{enumerate}

\hfill\break
\hfill\break
\hfill\break
\hfill\break
\hfill\break
\hfill\break
\hfill\break

\begin{enumerate}
\def\labelenumi{\arabic{enumi}.}
\setcounter{enumi}{7}
\tightlist
\item
  Is the original WHI study from 2002 then wrong? Why does it generate a
  statistically significant result for all ages?
\end{enumerate}

\hfill\break
\hfill\break
\hfill\break
\hfill\break
\hfill\break
\hfill\break
\hfill\break
\hfill\break

\begin{enumerate}
\def\labelenumi{\arabic{enumi}.}
\setcounter{enumi}{8}
\tightlist
\item
  In order to answer questions at a more subgroup level, do you think
  more power is required?
\end{enumerate}

\hfill\break
\hfill\break
\hfill\break
\hfill\break
\hfill\break
\hfill\break
\hfill\break

\subsection{Data Analysis}\label{data-analysis}

We shall now proceed to analyze some of these questions on two simulated
datasets. The datasets are stored in a \texttt{.csv} file. As earlier,
use the \texttt{readr} package, which is included inside the
\texttt{tidyverse}. If you haven't installed the tidyverse before, you
can do so by running \texttt{install.packages("tidyverse")} once.

\begin{Shaded}
\begin{Highlighting}[]
\CommentTok{\# load tidyverse (includes readr for CSVs)}
\FunctionTok{library}\NormalTok{(tidyverse)}
\end{Highlighting}
\end{Shaded}

\begin{Shaded}
\begin{Highlighting}[]
\CommentTok{\# load the data into R}

\NormalTok{data\_randomized }\OtherTok{\textless{}{-}} \FunctionTok{read\_csv}\NormalTok{(}\StringTok{"https://stat158.berkeley.edu/spring{-}2026/data/hormone{-}risk/hormone\_risk\_randomized.csv"}\NormalTok{)}

\NormalTok{data\_observations }\OtherTok{\textless{}{-}} \FunctionTok{read\_csv}\NormalTok{(}\StringTok{"https://stat158.berkeley.edu/spring{-}2026/data/hormone{-}risk/hormone\_risk\_observational.csv"}\NormalTok{)}
\end{Highlighting}
\end{Shaded}

The dataset contains three major columns: Age Group (50-59/ 60-69/
70-79), Treatment (1/0), and Response (Continuous). The Binary Treatment
here reflects a single form and dosage level of the hormonal supplement,
evaluated against a continuous health outcome indicating cardiovascular
risk. The observational dataset contains an additional baseline health
score variable.

\subsubsection{@. T-tests and Power
calculations}\label{t-tests-and-power-calculations}

First, proceed to conduct t-tests on the whole population and at the
age-group levels using the inbuilt \texttt{t.test()} function in
\texttt{R}. Do this for both the randomized and observational datasets.

\begin{Shaded}
\begin{Highlighting}[]
\CommentTok{\# Conduct t{-}tests}
\end{Highlighting}
\end{Shaded}

Next, investigate what the required sample size for the power of a
t-test would be here for the overall population. Use the observed effect
in the whole population for your calculations. Fill in the rest of the
\texttt{power.t.test()} function from \texttt{R} below. The data was
simulated using Gaussian distributions, so a t-test is valid in this
case.

\begin{Shaded}
\begin{Highlighting}[]
\CommentTok{\# Calculate Effect Size from Population. Fill in the following:}

\CommentTok{\# overall\_mean\_diff \textless{}{-} with(\# data,}
\CommentTok{\#                           mean(response[treatment==\# Fill In]) {-}}
\CommentTok{\#                             mean(response[treatment==\# Fill In]))}

\CommentTok{\# overall\_sd \textless{}{-} sd(data$response)}

\CommentTok{\# Check Sample size for 80\% Power}
\FunctionTok{power.t.test}\NormalTok{(}
  \AttributeTok{delta =}\NormalTok{ overall\_mean\_diff,}
  \AttributeTok{sd =}\NormalTok{ overall\_sd,}
  \AttributeTok{sig.level =} \FloatTok{0.05}\NormalTok{,}
  \AttributeTok{power =} \FloatTok{0.8}\NormalTok{,}
  \AttributeTok{type =} \StringTok{"two.sample"}
\NormalTok{)}
\end{Highlighting}
\end{Shaded}

\subsubsection{@. Investigating the Bias}\label{investigating-the-bias}

We first subsample randomly from the randomized study to reflect a
smaller RCT. We then perform the usual t-test to see if its significant.
Feel free to change the size of the smaller sample above and below your
previous threshold and see how that changes effects.

\begin{Shaded}
\begin{Highlighting}[]
\CommentTok{\# Subsampling}

\NormalTok{small\_study\_sim }\OtherTok{\textless{}{-}} \ControlFlowTok{function}\NormalTok{(}\AttributeTok{n\_sample =} \DecValTok{400}\NormalTok{) \{}
  
\NormalTok{  sample\_data }\OtherTok{\textless{}{-}}\NormalTok{ data\_randomized[}\FunctionTok{sample}\NormalTok{(}\DecValTok{1}\SpecialCharTok{:}\FunctionTok{nrow}\NormalTok{(data\_randomized), n\_sample), ]}
  
  \CommentTok{\# result \textless{}{-} t.test()}
  
  \FunctionTok{return}\NormalTok{(result}\SpecialCharTok{$}\NormalTok{p.value)}
\NormalTok{\}}
\end{Highlighting}
\end{Shaded}

Now, test how often this subsample turns out to be significant with
\texttt{alpha\ =\ 0.05}

\begin{Shaded}
\begin{Highlighting}[]
\CommentTok{\# Run multiple small studies}
\NormalTok{n\_sim }\OtherTok{\textless{}{-}} \DecValTok{1000}
\NormalTok{p\_values }\OtherTok{\textless{}{-}} \FunctionTok{replicate}\NormalTok{(n\_sim, }\FunctionTok{small\_study\_sim}\NormalTok{(}\DecValTok{400}\NormalTok{))}

\FunctionTok{mean}\NormalTok{(p\_values }\SpecialCharTok{\textless{}} \FloatTok{0.05}\NormalTok{)}
\end{Highlighting}
\end{Shaded}

We next simulate a method of sampling with bias from the observational
dataset. This is outlined below.

\begin{Shaded}
\begin{Highlighting}[]
\CommentTok{\# Subsampling with bias}

\NormalTok{small\_biased\_sim }\OtherTok{\textless{}{-}} \ControlFlowTok{function}\NormalTok{(}\AttributeTok{n\_sample =} \DecValTok{400}\NormalTok{) \{}
  
  \CommentTok{\# Only allow treated women from the top 40\% healthiest}
\NormalTok{  cutoff }\OtherTok{\textless{}{-}} \FunctionTok{quantile}\NormalTok{(data\_observational}\SpecialCharTok{$}\NormalTok{health, }\FloatTok{0.60}\NormalTok{)}
  
\NormalTok{  eligible }\OtherTok{\textless{}{-}} \FunctionTok{with}\NormalTok{(data\_observational,}
\NormalTok{                   (treatment }\SpecialCharTok{==} \StringTok{"Hormone"} \SpecialCharTok{\&}\NormalTok{ health }\SpecialCharTok{\textgreater{}}\NormalTok{ cutoff) }\SpecialCharTok{|}
\NormalTok{                     (treatment }\SpecialCharTok{==} \StringTok{"Control"}\NormalTok{))}
  
\NormalTok{  pool }\OtherTok{\textless{}{-}}\NormalTok{ data\_observational[eligible, ]}
  
\NormalTok{  sample\_data }\OtherTok{\textless{}{-}}\NormalTok{ pool[}\FunctionTok{sample}\NormalTok{(}\DecValTok{1}\SpecialCharTok{:}\FunctionTok{nrow}\NormalTok{(pool), n\_sample), ]}
  
  \CommentTok{\# result \textless{}{-} t.test(data = sample\_data)}
  
  \FunctionTok{c}\NormalTok{(}
    \AttributeTok{p =}\NormalTok{ result}\SpecialCharTok{$}\NormalTok{p.value,}
    \AttributeTok{effect =} \FunctionTok{as.numeric}\NormalTok{(result}\SpecialCharTok{$}\NormalTok{estimate[}\StringTok{"mean in group Hormone"}\NormalTok{] }\SpecialCharTok{{-}}
\NormalTok{                          result}\SpecialCharTok{$}\NormalTok{estimate[}\StringTok{"mean in group Control"}\NormalTok{])}
\NormalTok{  )}
\NormalTok{\}}
\end{Highlighting}
\end{Shaded}

Again conduct the t-test with these new subsample. How often are they
significant? And what is the direction?

\begin{Shaded}
\begin{Highlighting}[]
\CommentTok{\# Run multiple small studies}

\NormalTok{n\_sim }\OtherTok{\textless{}{-}} \DecValTok{1000}

\NormalTok{results\_biased }\OtherTok{\textless{}{-}} \FunctionTok{replicate}\NormalTok{(n\_sim, }\FunctionTok{small\_biased\_sim}\NormalTok{(}\DecValTok{400}\NormalTok{))}

\CommentTok{\# Proportion statistically significant}
\FunctionTok{mean}\NormalTok{(results\_biased[}\StringTok{"p"}\NormalTok{, ] }\SpecialCharTok{\textless{}} \FloatTok{0.05}\NormalTok{)}

\CommentTok{\# Average p{-}value}
\FunctionTok{mean}\NormalTok{(results\_biased[}\StringTok{"p"}\NormalTok{, ])}

\CommentTok{\# Proportion with protective effect (negative estimate)}
\FunctionTok{mean}\NormalTok{(results\_biased[}\StringTok{"effect"}\NormalTok{, ] }\SpecialCharTok{\textless{}} \DecValTok{0}\NormalTok{)}
\end{Highlighting}
\end{Shaded}

\subsubsection{@. Power Dilution with
Subgroups}\label{power-dilution-with-subgroups}

Now, calculate using \texttt{power.t.test()} the required sample size
for each age-group and the observed effect size in each group. Assume
independent two-sample t-tests in each age group.

\begin{Shaded}
\begin{Highlighting}[]
\CommentTok{\# Power calculations}

\DocumentationTok{\#\# 50{-}59}

\NormalTok{subgroup\_data }\OtherTok{\textless{}{-}} \FunctionTok{subset}\NormalTok{(data\_randomized, age\_group }\SpecialCharTok{==} \StringTok{"50{-}59"}\NormalTok{)}

\CommentTok{\# sub\_diff \textless{}{-} with(subgroup\_data,}
\CommentTok{\#                  mean(response[treatment==?]) {-}}
\CommentTok{\#                    mean(response[treatment==?]))}

\NormalTok{sub\_sd }\OtherTok{\textless{}{-}} \FunctionTok{sd}\NormalTok{(subgroup\_data}\SpecialCharTok{$}\NormalTok{response)}

\FunctionTok{power.t.test}\NormalTok{(}
  \AttributeTok{delta =}\NormalTok{ sub\_diff,}
  \AttributeTok{sd =}\NormalTok{ sub\_sd,}
  \AttributeTok{sig.level =} \FloatTok{0.05}\NormalTok{,}
  \AttributeTok{power =} \FloatTok{0.8}\NormalTok{,}
  \AttributeTok{type =} \StringTok{"two.sample"}\NormalTok{,}
  \AttributeTok{alternative =} \StringTok{"two.sided"}
\NormalTok{)}

\DocumentationTok{\#\# Repeat for other two subgroups}
\end{Highlighting}
\end{Shaded}

Sum these up over all three age groups. Is this larger or smaller than
the total sample size you calculated from Question 10?

\begin{Shaded}
\begin{Highlighting}[]
\CommentTok{\# Sum individual power calculations and compare}
\end{Highlighting}
\end{Shaded}

How does this compare to the 2002 and 2007 WHI studies and their
results?

\hfill\break
\hfill\break
\hfill\break
\hfill\break
\hfill\break




\end{document}
